\section{Ce que disent les magazines et la recherche}\label{sec:presse}
% Chapitre 9 (tache 10, passage chapitres 09-10). Presse: citations uniquement via le
% tableau genere \tabellenkoerper{data/tab_presse} (citations verifiees mot pour mot par
% scripts/presse.py contre refs/presse/<id>.md); la prose paraphrase, elle ne recopie
% aucune citation. Chaque id de data/presse.json est cite une seule fois par \QP dans ce
% chapitre. Comptes de la presse: _macros_redaction_kapitel_9_10 (rechnung_retraite.py),
% qui leve une erreur si les citations sur la fiscalite cessent d'etre toutes "contra".
% Recherche: data/literatur.bib reste vide (Scopus sans cle, Semantic Scholar HTTP 429 le
% 30.09.2026, deux essais); aucune reference citee de memoire (voir LUECKEN.md).
% Phrases liees aux donnees actuelles (a relire si presse.json change): "corpus domine par
% les arguments favorables", "les objections se concentrent sur l'impot et la phase de
% retrait", colonne "Accord" du tableau tab:accord-presse-document.

Ce chapitre r\'epond \`a la question~: que disent la presse sp\'ecialis\'ee allemande et la recherche
de la retraite par les dividendes, et les r\'esultats de ce document leur donnent-ils raison~? Le
chapitre~\ref{sec:retraite} s'est conclu sur trois constats~: pas de phase pont au sc\'enario de base,
une r\'eserve d'un an qui apporte l'essentiel de la protection mesurable, et une r\`egle des
\TauxRetrait~\% qui, rejou\'ee sur l'histoire am\'ericaine, tient plus souvent que le revenu de
dividendes avec moins de capital. La section~\ref{sec:presse-corpus} d\'ecrit les articles lus et la
fa\c{c}on dont leurs citations ont \'et\'e v\'erifi\'ees. Les sections~\ref{sec:presse-sans-vendre}
\`a~\ref{sec:presse-epargne} confrontent leurs arguments, th\`eme par th\`eme, aux chiffres des
chapitres pr\'ec\'edents, et la section~\ref{sec:presse-bilan} r\'esume ces accords et d\'esaccords
dans un tableau. La section~\ref{sec:presse-recherche} dit ce qu'il en est de la recherche
acad\'emique.

\subsection{Ce qui a \'et\'e lu, et comment}\label{sec:presse-corpus}

\annahme{Corpus de presse~: \PresseCitations{} citations tir\'ees de \PresseArticles{} articles de
\PressePublications{} publications allemandes, magazines financiers, sites sp\'ecialis\'es et blogs de
la communaut\'e FIRE (ind\'ependance financi\`ere, retraite anticip\'ee), consult\'es le
\PresseDateConsultation. Une citation n'est retenue que si elle figure mot pour mot dans le texte
t\'el\'echarg\'e de la page~; le script qui le v\'erifie \'ecarte toute citation introuvable. La
position de chaque citation (\enquote{pro}, \enquote{contra} ou \enquote{neutre} envers la strat\'egie
des dividendes) et son th\`eme ont \'et\'e attribu\'es \`a la collecte, avant la r\'edaction de ce
chapitre. Ce sont des sources secondaires~: elles expriment une opinion, et ce plan n'en reprend
aucune valeur dans ses calculs.}

Le tableau~\ref{tab:citations-presse-verifiees} reproduit ces citations dans leur langue d'origine,
abr\'eg\'ees par des points de suspension quand elles sont longues~; chaque article est r\'ef\'erenc\'e
en note \`a sa premi\`ere mention dans ce chapitre, et en annexe. Le corpus compte \PressePro{}
citations favorables, \PresseContra{} d\'efavorables et \PresseNeutre{} neutre. Ce d\'ecompte n'est pas
un vote de la presse allemande~: il refl\`ete ce qui \'etait lisible la nuit de la collecte. Deux
sources pr\'evues, Finanztest et le Handelsblatt, n'ont renvoy\'e qu'une page de navigation ou un court
chapeau derri\`ere leur acc\`es payant, et un troisi\`eme magazine n'\'etait pas accessible \`a l'outil de
recherche~; un second blog FIRE les remplace en partie. Les \'etiquettes, enfin, sont discutables par
endroits. La phrase de Finanztip class\'ee \enquote{pro} sous le th\`eme des ETF distribuants
recommande en r\'ealit\'e les ETF capitalisants pour constituer un patrimoine, ce qui d\'ecrit la
r\'ef\'erence de ce plan plut\^ot que la strat\'egie des dividendes, et la phrase de B\"orse Online sur
les actions \`a dividendes porte le m\^eme th\`eme sans parler d'ETF. \bk{Ce chapitre juge donc chaque
citation sur son texte, pas sur son \'etiquette.}

{\footnotesize
\renewcommand{\_}{\textunderscore\allowbreak}
\setlength{\LTpre}{\medskipamount}\setlength{\LTpost}{\medskipamount}
\begin{longtable}{>{\bfseries\raggedright\arraybackslash}p{0.22\linewidth}
  >{\raggedright\arraybackslash}p{0.08\linewidth} >{\raggedright\arraybackslash}p{0.15\linewidth}
  >{\raggedright\arraybackslash}p{0.44\linewidth}}
  \caption{Citations de presse v\'erifi\'ees mot pour mot dans le texte t\'el\'echarg\'e des articles,
  avec la position et le th\`eme attribu\'es \`a la collecte (sources secondaires, consult\'ees le
  \PresseDateConsultation).}\label{tab:citations-presse-verifiees}\\
  \toprule
  \rowcolor[gray]{0.9}
  \textbf{Source} & \textbf{Position} & \textbf{Th\`eme (code)} & \textbf{Citation (extrait)} \\
  \midrule
  \endfirsthead
  \toprule
  \rowcolor[gray]{0.9}
  \textbf{Source} & \textbf{Position} & \textbf{Th\`eme (code)} & \textbf{Citation (extrait)} \\
  \midrule
  \endhead
  \bottomrule
  \endlastfoot
  \tabellenkoerper{data/tab_presse}
\end{longtable}}

Ce tableau montre un corpus domin\'e par les arguments favorables, \PressePro{} citations sur
\PresseCitations, mais o\`u les objections se concentrent sur deux sujets pr\'ecis~: l'imp\^ot et la
phase de retrait. Les sections suivantes reprennent ces arguments un par un.

\subsection{Vivre sans rien vendre}\label{sec:presse-sans-vendre}

L'argument le plus fr\'equent du corpus est aussi le plus simple~: les dividendes arrivent sur le
compte sans qu'il faille vendre de parts, et surtout pas en pleine baisse des cours. B\"orse Online
le pr\'esente comme l'avantage technique des actions \`a dividendes\QP{boerseonline_fire_portfolio},
getmad.de comme un moyen de couvrir une partie des d\'epenses de la phase de retrait\QP{getmad_fire},
et justETF y ajoute deux effets~: il serait plus facile de garder son calme en crise, et le risque de
s\'equence comme le risque de faillite seraient r\'eduits au minimum, puisque le capital reste
investi\QP{justetf_entnahmestrategien}.

Ce document donne raison \`a la moiti\'e de cet argument. Le capital n'est jamais vendu, il ne peut
donc pas s'\'epuiser. Mais le revenu suit le dividende~: celui de l'indice \IndiceActionsUs{} a perdu
\GroessterEinbruch~\% de sa valeur r\'eelle de \GroessterEinbruchDebut{} \`a \GroessterEinbruchJahr, et
encore \EinbruchRecent~\% de \EinbruchRecentDebut{} \`a \EinbruchRecentFin{}
(section~\ref{sec:risques-baisses}). Un choc de $-$20~points sur la croissance du dividende la
premi\`ere ann\'ee de retraite fait tomber la r\'eussite du mod\`ele sans imp\^ot de
\ErfolgsquoteBasis{} \`a \McErfolgChocMoinsVingt~\% (section~\ref{sec:risques-sequence})~: le risque
de s\'equence ne dispara\^it pas, il passe du capital au revenu. Sur les m\^emes d\'eparts historiques,
enfin, la r\`egle des \TauxRetrait~\%, qui vend des parts chaque ann\'ee, tient dans
\RetraitHistMemesDepartsProzent~\% des cas, contre \RueckspielTenusProzent~\% pour le revenu de
dividendes (section~\ref{sec:retraite-quatre}).

L'effet psychologique, en revanche, est r\'eel, et ce document ne le mesure pas. Ne pas avoir \`a
d\'ecider de vendre au pire moment a de la valeur pour qui craint de paniquer. FinanzNachrichten.de
r\'esume le d\'esaccord en une phrase~: la strat\'egie des dividendes serait un outil psychologique
plut\^ot qu'une strat\'egie math\'ematiquement sup\'erieure\QP{finanznachrichten_irrweg}. Les chiffres
de ce plan vont dans ce sens~; \bk{c'est \`a toi de dire combien ce confort vaut pour toi, en capital
et en ann\'ees de travail.}

\subsection{Le prix de ce confort~: imp\^ot, capital et \^age}\label{sec:presse-prix}

Les \PresseCitationsFiscalite{} citations du corpus qui portent sur l'imp\^ot sont toutes
d\'efavorables \`a la strat\'egie des dividendes, et c'est le point o\`u ce document leur donne le plus
nettement raison. justETF rappelle que tout le dividende vers\'e subit l'imp\^ot sur les revenus du
capital, la contribution de solidarit\'e et, le cas \'ech\'eant, l'imp\^ot d'\'Eglise~; getmad.de, que
cet imp\^ot abaisse le taux de retrait effectif~; B\"orse Online range la fiscalit\'e allemande parmi
les obstacles \`a la constitution d'un patrimoine, \`a c\^ot\'e du co\^ut de la vie et des
loyers\QP{boerseonline_fire_traum}.

Le chapitre~\ref{sec:fiscalite} chiffre cet \'ecart~: sur cent euros de dividende brut, il te reste
\NettoDe~€ pour une action allemande et \NettoUs~€ pour une action am\'ericaine, contre \NettoEtf~€
pour un ETF actions, et l'assurance maladie volontaire prend en plus \KvPvAnteilBrut~\% du revenu brut
dont tu as besoin. Le chapitre~\ref{sec:capital} en tire la cons\'equence~: \`a revenu net \'egal, la
r\'epartition \ReferenzStrategie{} demande \KapitalMaisonSiebzig~€ de capital, soit
\KapitalDifferenzSiebzigProzent~\% de plus que les \KapitalEtfReferenz~€ d'un ETF monde avec retrait de
\TauxRetrait~\%. Le chapitre~\ref{sec:accumulation} traduit cet \'ecart en ann\'ees~: au taux
d'\'epargne de base, la r\'ef\'erence atteint l'objectif \EcartAgeEtfMaisonBasis~ans plus t\^ot.

FinanzNachrichten.de ajoute que le confort mental se paie en diversification et en rendement, et
d\'econseille l'approche pour la phase de retrait. Sur la diversification, le
chapitre~\ref{sec:portefeuille} lui donne en partie raison~: un ETF maison de \NombreTitres{} titres
supprime l'essentiel du risque propre \`a chaque titre, mais il ne couvre que \PortefeuillePays{} pays,
et rien ne montre que sa m\'ethode bat un ETF monde. Sur la phase de retrait, le test historique de la
section~\ref{sec:retraite-quatre} va dans le m\^eme sens, avec une nuance d\'ej\`a dite~: un revenu de
dividendes qui baisse reste un revenu, alors qu'un capital \'epuis\'e ne verse plus rien.

\subsection{Combien retirer chaque ann\'ee~?}\label{sec:presse-retrait}

Deux articles parlent du taux de retrait, avec des conseils oppos\'es. B\"orse Online pr\'esente la
r\`egle des quatre pour cent comme une formule empirique, prudente au regard de rendements historiques
avant imp\^ot de huit pour cent et plus sur les grands indices d'actions. Finanztip conseille au
contraire de ne retirer chaque ann\'ee qu'un \`a trois pour cent du capital r\'euni au d\'epart en
retraite\QP{finanztip_geldanlage}.

Les donn\'ees am\'ericaines de ce document donnent plut\^ot raison au second. Un retrait de
\TauxRetrait~\% du capital initial, ajust\'e de l'inflation, a tenu \RetraitHistAnsVierzig~ans dans
\RetraitHistReussiteVierzig~\% des d\'eparts depuis \RetraitHistPremierVierzig, mais il a \'epuis\'e le
capital des retraites commenc\'ees en \RetraitHistEchecsListeVierzig~; \`a \TauxRetraitPrudent~\%,
aucun d\'epart n'a \'echou\'e (section~\ref{sec:retraite-quatre}). La formule n'est donc pas prudente
dans tous les cas, et le chiffre de huit pour cent qu'invoque B\"orse Online n'est pas comparable \`a
ceux de ce plan~: c'est un rendement avant imp\^ot et avant inflation, alors que ce document raisonne en
rendement r\'eel, \RenditeReal~\% par an. Ce taux est l'un des deux \'ecarts assum\'es au plan~: c'est
le \RenditeReelleConvention{} des rendements annuels, parce que seule la moyenne g\'eom\'etrique d\'ecrit
ce que devient un capital apr\`es des d\'ecennies de hausses et de krachs (chapitre~\ref{sec:depart}).
Le test historique ignore de plus l'imp\^ot allemand et les frais, qui ne peuvent qu'abaisser ces taux
de r\'eussite.

La prudence de Finanztip a un prix~: \`a \TauxRetraitPrudent~\%, l'ETF monde demande
\KapitalEtfPrudent~€ au lieu de \KapitalEtfReferenz~€. \bk{Elle reste pourtant moins ch\`ere que la
strat\'egie des dividendes seuls, qui demande \KapitalMaisonSiebzig~€ dans la r\'epartition de base.}

\subsection{Qualit\'e, diversification et secteurs}\label{sec:presse-selection}

Sur le choix des titres, la presse est plus nuanc\'ee que sur la retraite elle-m\^eme. B\"orse Online
d\'econseille de se fier au seul rendement du dividende et recommande de fixer d'abord un m\'elange
d'actions, d'ETF et \'eventuellement d'obligations, puis de chercher la qualit\'e, des mod\`eles
d'affaires conservateurs et de vraies valeurs de rendement. justETF signale, pour les investisseurs
prudents, les entreprises qui ont maintenu ou relev\'e leur dividende pendant de nombreuses ann\'ees,
les aristocrates du dividende, et rappelle qu'un ETF distribuant dispense de choisir soi-m\^eme les
titres\QP{justetf_ausschuettend}.

La m\'ethode du chapitre~\ref{sec:portefeuille} suit ces conseils~: elle borne le rendement entre
\RegleRendementMin{} et \RegleRendementMax~\%, limite le taux de distribution \`a \ReglePayoutMax~\% du
b\'en\'efice, exige un dividende sans baisse sur \RegleAnneesCroissance~ans, puise une partie de son
univers dans les aristocrates du dividende am\'ericains et garde une part d'ETF monde. Le m\^eme
chapitre en montre les limites~: le crit\`ere de r\'egularit\'e regarde en arri\`ere et n'aurait pas
\'ecart\'e \BeispielBaisse{} la veille de la suppression de son dividende en \BeispielBaisseAnnee, un
univers d'aristocrates porte par construction un biais du survivant, et aucun test historique ne
montre que la m\'ethode fait mieux qu'un ETF monde (section~\ref{sec:porte-biais}).

Le blog Dividende statt Rente affirme que les entreprises \`a haut dividende fluctuent moins et
r\'esistent mieux aux crises\QP{dividende_statt_rente}. Ce document ne compare pas la fluctuation de
leurs cours \`a celle du march\'e et ne peut ni confirmer ni infirmer cette premi\`ere moiti\'e de la
phrase~; il montre en revanche que leur dividende, lui, n'est pas \`a l'abri des crises
(section~\ref{sec:risques-baisses}). Extra-Magazin observe enfin que, dans le DAX, ce sont surtout les
valeurs financi\`eres qui augmentent nettement leurs distributions\QP{extraetf_dividendenboom}. C'est
le constat d'une p\'eriode, pas une r\`egle de construction~: la m\'ethode du
chapitre~\ref{sec:portefeuille} limite chaque secteur \`a \RegleSecteurMax~titres, pour qu'aucun ne
domine l'ETF maison, quelle que soit la mode du moment.

\subsection{\'Epargner, investir en actions, battre l'inflation}\label{sec:presse-epargne}

Le dernier groupe de citations porte moins sur les dividendes que sur l'\'epargne elle-m\^eme, et c'est
l\`a que la presse et ce document s'accordent le mieux. Extra-Magazin rappelle que l'on n'atteint pas
les objectifs du mouvement FIRE en laissant son argent sur un compte d'\'epargne dont les int\'er\^ets
ne compensent pas l'inflation\QP{extraetf_fire}. Finanztip juge une forte part d'actions importante
pour la retraite, les ETF actions adapt\'es comme source de rendement au-dessus de l'inflation, et
les ETF capitalisants particuli\`erement adapt\'es \`a la constitution d'un
patrimoine\QP{finanztip_thesaurierend}~; cette derni\`ere phrase d\'ecrit exactement la r\'ef\'erence de
ce plan. L'inflation n'a rien d'abstrait~: la section~\ref{sec:risques-inflation} a montr\'e ce que
la pouss\'ee de \InflationPicAnnee{} fait \`a un revenu qui n'est pas index\'e.

Deux arguments favorables aux dividendes restent \`a nuancer. Dividende statt Rente pr\'esente les
dividendes comme une alternative aux obligations tant que les taux d'int\'er\^et restent bas~; cet
argument d\'epend de l'\'epoque o\`u il a \'et\'e \'ecrit, et ce plan, qui ne contient aucune obligation,
ne permet pas de le v\'erifier. justETF affirme qu'investir en actions et en obligations en encaissant
dividendes et int\'er\^ets peut te faire atteindre ton objectif. Ce document confirme que c'est
possible, mais pas t\^ot~: au taux de base, la strat\'egie des dividendes atteint l'objectif \`a
\AlterYannBasis~ans, l'\^age l\'egal, et \`a \SparquoteZwanzig~\% d'\'epargne elle ne l'atteint jamais
(section~\ref{sec:accu-age}). \bk{Aucun article du corpus ne chiffre un \^age de d\'epart~; c'est
pourtant la seule question qui compte pour toi.}

\subsection{Accords et d\'esaccords}\label{sec:presse-bilan}

Le tableau~\ref{tab:accord-presse-document} met en regard, th\`eme par th\`eme, ce qu'affirme le
corpus et ce que montrent les calculs de ce document. La derni\`ere colonne ne juge pas la qualit\'e des
articles~; elle dit seulement si les chiffres de ce plan vont dans leur sens.

\begin{table}[!tb]
  \centering\small
  \caption{Arguments du corpus de presse face aux r\'esultats de ce document (r\`egles de 2026, euros
  de 2026, sc\'enario de base \ReferenzStrategie{} \`a \ReferenzSparquote~\% d'\'epargne).}
  \label{tab:accord-presse-document}
  \begin{tabularx}{\linewidth}{>{\bfseries\raggedright\arraybackslash}p{0.16\linewidth}
    >{\raggedright\arraybackslash}X >{\raggedright\arraybackslash}X
    >{\raggedright\arraybackslash}p{0.13\linewidth}}
    \toprule
    \rowcolor[gray]{0.9}
    \textbf{Th\`eme} & \textbf{Ce qu'affirme le corpus} & \textbf{Ce que montre ce document}
      & \textbf{Accord} \\
    \midrule
    Vivre sans vendre & Capital intact, calme en crise, risque de s\'equence minimal
      & Capital intact, mais revenu expos\'e~: dividende r\'eel de l'indice en baisse de
      \GroessterEinbruch~\% (section~\ref{sec:risques-baisses}) & En partie \\
    Imp\^ot & Tout le dividende est impos\'e (\PresseCitationsFiscalite{} citations, toutes
      d\'efavorables) & \NettoEtf~€ nets pour cent euros d'ETF, \NettoDe~€ pour une action
      allemande (chapitre~\ref{sec:fiscalite}) & Oui \\
    Capital et \^age & Confort pay\'e en rendement et en diversification
      & \KapitalDifferenzSiebzigProzent~\% de capital en plus, objectif atteint
      \EcartAgeEtfMaisonBasis~ans plus tard (chapitres~\ref{sec:capital}
      et~\ref{sec:accumulation}) & Oui \\
    Taux de retrait & \TauxRetrait~\% prudent (B\"orse Online), un \`a trois pour cent (Finanztip)
      & \TauxRetrait~\%~: \RetraitHistReussiteVierzig~\% de r\'eussite sur \RetraitHistAnsVierzig~ans~;
      \TauxRetraitPrudent~\%~: \RetraitHistReussitePrudentVierzig~\% (section~\ref{sec:retraite-quatre})
      & Plut\^ot Finanztip \\
    Choix des titres & Qualit\'e, aristocrates, pas le seul rendement
      & R\`egles conformes, mais biais du survivant et aucun test historique
      (section~\ref{sec:porte-biais}) & En partie \\
    Stabilit\'e en crise & Moins de fluctuations, meilleure r\'esistance
      & Fluctuation des cours non mesur\'ee~; dividendes coup\'es en crise
      (section~\ref{sec:risques-baisses}) & Non d\'emontr\'e \\
    Actions et inflation & L'\'epargne sur compte ne suffit pas, forte part d'actions
      & Plan fond\'e sur un rendement r\'eel de \RenditeReal~\% par an en actions
      (chapitre~\ref{sec:depart}) & Oui \\
    \bottomrule
  \end{tabularx}
\end{table}

Ce tableau montre que ce document rejoint la presse sur l'\'epargne, les actions et l'imp\^ot, et
qu'il s'en \'ecarte sur ce qu'elle pr\'esente comme la force des dividendes~: un revenu stable et
s\^ur. Les d\'esaccords portent tous sur ce m\^eme point, que les chapitres~\ref{sec:risques}
et~\ref{sec:retraite} montrent moins solide que ne le dit le corpus.

\subsection{Et la recherche~?}\label{sec:presse-recherche}

Les \'etudes acad\'emiques pr\'evues pour ce chapitre, sur les taux de retrait soutenables, le risque
de s\'equence, la diversification selon le nombre de titres, la performance des strat\'egies \`a haut
dividende et la pr\'ef\'erence des investisseurs pour les dividendes, n'ont pas pu \^etre v\'erifi\'ees,
la base Scopus \'etant inaccessible et l'interface publique de Semantic Scholar ayant refus\'e les
requ\^etes lors de deux essais le 30~septembre~2026, si bien que ce chapitre ne cite aucune \'etude
plut\^ot que d'en citer de m\'emoire.

Une partie de ces questions trouve malgr\'e tout une r\'eponse dans les calculs de ce document, avec
leurs limites. Le taux de retrait et le risque de s\'equence sont test\'es sur l'histoire de
l'\IndiceActionsUs{} depuis \ShillerDebut{} (chapitres~\ref{sec:risques} et~\ref{sec:retraite}), mais
sur ce seul march\'e et sans imp\^ot. La diversification selon le nombre de titres est mesur\'ee sur les
\NombreTitres{} titres du portefeuille-exemple et sur \RegleHistoriqueAns~ans de cours seulement
(section~\ref{sec:porte-diversification}). La performance compar\'ee des strat\'egies \`a haut dividende
et la pr\'ef\'erence des investisseurs pour les dividendes ne sont pas trait\'ees. \bk{Confronter ces
r\'esultats \`a des \'etudes publi\'ees et v\'erifi\'ees reste \`a faire~; c'est la lacune la plus
importante de ce chapitre.}

\bk{Retiens de ce chapitre que la presse vante surtout le confort de vivre sans rien vendre, et que
ce document confirme ce confort mais en chiffre le prix~: \KapitalDifferenzSiebzigProzent~\% de
capital en plus, \EcartAgeEtfMaisonBasis~ans de travail en plus au taux de base, et un revenu qui
baisse avec le dividende dans les crises.} Sur l'imp\^ot et sur l'int\'er\^et d'investir en actions, la
presse et ce document disent la m\^eme chose. Le chapitre~\ref{sec:plan} traduit l'ensemble de ces
r\'esultats en un plan d'action dat\'e.

\Yannbox{\AnneeYannBasis}{Si Yann suit l'argument le plus r\'epandu de la presse, vivre des dividendes
sans rien vendre, il doit r\'eunir \KapitalMaisonSiebzig~€ avec la r\'epartition \ReferenzStrategie{}
et ne s'arr\^ete qu'en \AnneeYannBasis, \`a \AlterYannBasis~ans. S'il suit les critiques, un ETF monde
avec un retrait de \TauxRetrait~\% par an, \KapitalEtfReferenz~€ suffisent, r\'eunis d\`es
\AnneeEtfBasis, \`a \AlterEtfBasis~ans, avec la m\^eme \'epargne. M\^eme avec la prudence de Finanztip,
\TauxRetraitPrudent~\% de retrait, il lui faudrait \KapitalEtfPrudent~€, moins que pour les dividendes
seuls. Sur ses dividendes, il paierait l'imp\^ot sur tout le montant vers\'e~: pour cent euros bruts
d'une action allemande, il garderait \NettoDe~€, contre \NettoEtf~€ pour un ETF actions.}
